\section{结论}\label{sec:conclusion}

本文综述了孪生素数猜想从1849年至今的研究历程。可以概括出三条主线：

\begin{enumerate}[label=(\arabic*)]
    \item \textbf{筛法主线}：从布伦的纯筛法、塞尔伯格的二次型筛法到陈景润的加权筛，筛法理论在“逼近到奇偶性壁垒”的意义上已臻完备；布朗定理与陈氏定理是这条路线上的两座丰碑。
    \item \textbf{间隔主线}：GPY 方法与张益唐的分布水平突破、Maynard--Tao 多维筛法一道，将有界素数间隔从猜想变为定理，间隔界从 $7 \times 10^7$ 压缩到 $246$，并贡献了素数分布研究近十年来最活跃的方向。
    \item \textbf{启发式主线}：Hardy--Littlewood 素数 $k$ 元组猜想与 Cram\'er 模型为孪生素数给出了量级正确的定量预言，计算验证（至 $10^{19}$）持续支持这些公式。
\end{enumerate}

孪生素数猜想与哥德巴赫猜想共享“奇偶性壁垒”这一根本障碍，与黎曼猜想共享“素数分布的深层随机性”这一核心谜题。张益唐2013年的突破证明了一个重要的事实：被认为“此生无望”的问题可以在现有工具的精妙组合下取得本质进展。未来的突破更可能来自观念的更新——如何在一个筛法框架内“看见”素因子个数的奇偶性——而非现有方法的继续打磨。
